\chapter*{Lời Mở Đầu}
\addcontentsline{toc}{chapter}{Lời Mở Đầu}
\markboth{Lời Mở Đầu}{Lời Mở Đầu}

\begin{truyen}{Lời Mở Đầu}{Opening Words}
\chuhoa{N}{hững câu chuyện hay nhất về học sinh không phải là những câu chuyện ngọt ngào.}
\textit{(The best stories about students are not the sweet ones.)}

Chúng là những câu chuyện thô ráp, trực tiếp và đôi khi khó nghe. Về điểm số, về áp lực, về so sánh, về bắt nạt, về mạng xã hội, về việc chọn nghề và về nỗi sợ bị thất bại.
\textit{(They are the rough, direct, and sometimes hard-to-hear ones. About grades, pressure, comparison, bullying, social media, career choices, and the fear of failure.)}

Quyển tám không né tránh những chuyện đó. Nó gọi thẳng tên những thứ học sinh đang sống qua mỗi ngày — bằng giọng thẳng, ngắn và thật.
\textit{(Volume eight does not avoid those things. It names directly what students are living through each day --- in a direct, brief, and honest voice.)}

Đọc quyển này không phải để cảm thấy tốt hơn ngay lập tức — mà để hiểu rõ hơn về thực tế mà mình đang đứng trong đó.
\textit{(Reading this volume is not to feel better immediately --- but to understand more clearly the reality in which you are standing.)}
\end{truyen}

\begin{baihoc}
Sự thật không phải lúc nào cũng dễ nghe --- nhưng nó là thứ duy nhất giúp ta thực sự thay đổi.
\textit{(The truth is not always easy to hear --- but it is the only thing that truly helps us change.)}
\end{baihoc}

\begin{ghinhoanh}
\textbf{Short English to remember:}

Honest stories help more than polished ones.

Students need truth, not only encouragement.
\end{ghinhoanh}

\begin{tuhocnhanh}
1. Trong những khó khăn học sinh thường gặp, điều nào em thấy mình đang trải qua? \textit{Among the common student struggles, which do you find yourself going through?}

2. Thẳng thắn với bản thân khó ở điểm nào? \textit{What makes it difficult to be honest with yourself?}

3. Em có thể làm gì để đối mặt với một áp lực đang tồn tại trong cuộc sống học đường của mình? \textit{What can you do to face a pressure currently existing in your school life?}
\end{tuhocnhanh}

\ngancach
